% =====================================================================
% v2 Researcher-Grade Extensions for ctexart / v2_minimal_report.tex
%
% Copy the relevant blocks into the main v2 report .tex when researcher-
% grade mode is on. The preamble macros below should be appended to the
% main report preamble (after the existing tcolorbox load) so the four
% colored boxes are available.
%
% Reference style: ProbeMate_CalibratedDiagnosticCommitment_CHI_CSCW_rewrite.tex
% =====================================================================

% --- PREAMBLE ADDITIONS ----------------------------------------------
% Add to v2_minimal_report.tex preamble if not already present:
%
%   \usepackage[most]{tcolorbox}
%   \usepackage{tabularx}
%   \usepackage{longtable}
%   \usepackage{booktabs}
%   \usepackage{array}
%   \usepackage{amsmath, amssymb}
%
% Then add the four colored box environments below.

\newtcolorbox{thesisbox}{
  colback=black!3,
  colframe=black!70,
  boxrule=0.6pt,
  arc=1.5mm,
  left=1.1em,
  right=1.1em,
  top=0.75em,
  bottom=0.75em
}

\newtcolorbox{bluebox}{
  colback=blue!4,
  colframe=blue!55!black,
  boxrule=0.5pt,
  arc=1.5mm,
  left=1.0em,
  right=1.0em,
  top=0.65em,
  bottom=0.65em
}

\newtcolorbox{redbox}{
  colback=red!3,
  colframe=red!60!black,
  boxrule=0.5pt,
  arc=1.5mm,
  left=1.0em,
  right=1.0em,
  top=0.65em,
  bottom=0.65em
}

\newtcolorbox{greenbox}{
  colback=green!4,
  colframe=green!45!black,
  boxrule=0.5pt,
  arc=1.5mm,
  left=1.0em,
  right=1.0em,
  top=0.65em,
  bottom=0.65em
}

% --- BODY BLOCKS (copy to body where indicated by integration map) ----

% =====================================================================
% Block 1 — Opening Thesis Box (insert before \section{一句话概括})
% =====================================================================

\begin{thesisbox}
\textbf{这版的核心决定：}TODO（3--6 句话，说明 \emph{核心构念决定}，不是项目主题）。
真正的构念不是 TODO，而是 TODO。因此本文把 \textit{[construct]} 定义为
TODO：\textbf{[dim 1]}、\textbf{[dim 2]}、\textbf{[dim 3]}。
\textbf{[Named levels]} 只是这个空间中的常见 archetype，不是真值类别、不是一维序列，
也不是专家必须一致标注的类别。
\end{thesisbox}

% =====================================================================
% Block 2 — Headline Competition Audit (subsection of 项目核心判断)
% =====================================================================

\subsection{标题候选与抢戏审计}

\begin{center}
\small
\begin{tabularx}{\textwidth}{>{\bfseries}p{0.22\textwidth}X X}
\toprule
候选 headline & 一句话内容 & 为什么不选作 spine \\
\midrule
TODO & TODO & TODO \\
TODO & TODO & TODO \\
\textbf{[Selected]} & TODO & \textbf{入选：TODO} \\
\bottomrule
\end{tabularx}
\end{center}

% =====================================================================
% Block 3 — Reverse Pitch (subsection of 核心问题不是)
% =====================================================================

\begin{quote}
本项目不是 TODO（misreading 1）；不是 TODO（misreading 2）；不是 TODO（misreading 3）。
它研究的是 TODO（precise object）。
\end{quote}

% =====================================================================
% Block 4 — 3D Construct Operationalization
% =====================================================================

\subsection{构念定义与操作化}

\begin{thesisbox}
\textbf{Definition.} \textit{[Construct name]} refers to the design problem of TODO.
\end{thesisbox}

\begin{center}
\small
\begin{tabularx}{\textwidth}{>{\bfseries}p{0.19\textwidth}p{0.16\textwidth}p{0.14\textwidth}p{0.14\textwidth}X}
\toprule
维度 & 0 & 1 & 2 & 3 \\
\midrule
Dim A\newline TODO & TODO & TODO & TODO & TODO \\
Dim B\newline TODO & TODO & TODO & TODO & TODO \\
Dim C\newline TODO & TODO & TODO & TODO & TODO \\
\bottomrule
\end{tabularx}
\end{center}

形式化地：
\[
K_t = (A_t, B_t, C_t)
\]
其中 $A_t$ 是 TODO，$B_t$ 是 TODO，$C_t$ 是 TODO。

% =====================================================================
% Block 5 — Archetype Table
% =====================================================================

\subsection{常见 archetype}

\begin{center}
\small
\begin{tabularx}{\textwidth}{>{\bfseries}p{0.16\textwidth}p{0.20\textwidth}p{0.24\textwidth}X}
\toprule
Archetype & Vector & 界面/系统形态 & 适用情境 \\
\midrule
TODO & $(0, 0, 0)$ & TODO & TODO \\
TODO & $(1, 2, 1)$ & TODO & TODO \\
TODO & $(2, 2, 2)$ & TODO & TODO \\
TODO & $(3, 3, 3)$ & TODO & TODO \\
\bottomrule
\end{tabularx}
\end{center}

\textbf{注意：}archetype 不是一维 ladder，也不是 ground-truth class。它是 design space
中常见落点的命名。

% =====================================================================
% Block 6 — Neighbor Concept Differentiation Table
% =====================================================================

\subsection{与邻居概念的区分}

\begin{center}
\small
\begin{tabularx}{\textwidth}{>{\bfseries}p{0.24\textwidth}X X}
\toprule
邻居概念 & 它主要解释什么 & 本构念的不同点 \\
\midrule
Model confidence & TODO & TODO \\
Selective prediction / abstention & TODO & TODO \\
Mixed-initiative interaction & TODO & TODO \\
Proactive vs reactive AI & TODO & TODO \\
Progressive disclosure & TODO & TODO \\
Appropriate reliance & TODO & TODO \\
\bottomrule
\end{tabularx}
\end{center}

具体差异化举例：本构念能描述 TODO（specific configuration），而 [neighbor X] 不能。

% =====================================================================
% Block 7 — Venue Strategy Table
% =====================================================================

\subsection{投稿策略}

\begin{bluebox}
\textbf{建议：}以 \textbf{TODO (primary venue)} 为主路线，但从第一天开始让数据也能支撑 \textbf{TODO (pivot venue)}。
TODO（pivot 触发条件）。
\end{bluebox}

\begin{center}
\small
\begin{tabularx}{\textwidth}{>{\bfseries}p{0.18\textwidth}X X}
\toprule
路线 & 主贡献写法 & 研究重心 \\
\midrule
TODO (Primary) & TODO & TODO \\
TODO (Pivot) & TODO & TODO \\
\bottomrule
\end{tabularx}
\end{center}

% =====================================================================
% Block 8 — RQ Matrix
% =====================================================================

\subsection{研究问题矩阵}

\begin{center}
\small
\begin{tabularx}{\textwidth}{>{\bfseries}p{0.06\textwidth}X p{0.12\textwidth} p{0.16\textwidth} p{0.18\textwidth} p{0.16\textwidth}}
\toprule
RQ & 研究问题 & Study & 主要 measure & 证据目标 & Fail condition \\
\midrule
RQ1 & TODO & Study 1 & TODO & TODO & TODO \\
RQ2 & TODO & Study 2 & TODO & TODO & TODO \\
RQ3 & TODO & Study 3 & TODO & TODO & TODO \\
RQ4 & TODO & Study 4 & TODO & TODO & TODO \\
RQ5 & TODO & Study 4 & TODO & TODO & TODO \\
\bottomrule
\end{tabularx}
\end{center}

% =====================================================================
% Block 9 — Alternative Outcomes Tables (Study 1)
% =====================================================================

\subsection{Study 1：形成性研究的可能发现}

不写"预期发现"，写"可能推翻或重塑设计的发现"：

\begin{center}
\small
\begin{tabularx}{\textwidth}{>{\bfseries}p{0.24\textwidth}X X}
\toprule
可能发现 & 它如何推翻/重塑原设计 & 设计后果 \\
\midrule
TODO (strong-form "design is wrong") & TODO & TODO \\
TODO (different-population finding) & TODO & TODO \\
TODO (novice vs expert inversion) & TODO & TODO \\
TODO (cost-too-high finding) & TODO & TODO \\
TODO ($\geq$ 5 rows for deep; $\geq$ 3 for standard) & TODO & TODO \\
\bottomrule
\end{tabularx}
\end{center}

% =====================================================================
% Block 10 — Deployment Cost as Finding + Two Scale Tiers
% =====================================================================

\subsection{部署的代价与两档规模}

\begin{redbox}
\textbf{必须承认：}TODO（deployment mechanism）不是免费的。每次 TODO（unit of cost）
需要 TODO 时间/工作量/设备。一节 TODO 分钟课做 TODO 次，占据约 TODO\% 的课堂/工作时间。
这不是"无缝叠加"，而是"有意插入的诊断节奏"——本身就是研究发现。
\end{redbox}

\begin{center}
\small
\begin{tabularx}{\textwidth}{>{\bfseries}p{0.26\textwidth}X}
\toprule
Trade-off & 研究意义 \\
\midrule
TODO & TODO \\
TODO & TODO \\
TODO & TODO \\
\bottomrule
\end{tabularx}
\end{center}

\subsubsection*{两档可行规模}

\begin{center}
\small
\begin{tabularx}{\textwidth}{>{\bfseries}p{0.18\textwidth}X X}
\toprule
方案 & 规模 & 适合路线 \\
\midrule
TODO ([primary]-强评估) & TODO & TODO \\
TODO ([pivot]-厚描) & TODO & TODO \\
\bottomrule
\end{tabularx}
\end{center}

% =====================================================================
% Block 11 — Reviewer Attack Table
% =====================================================================

\section{预期审稿人攻击与预防}

\begin{center}
\small
\begin{longtable}{>{\bfseries}p{0.24\textwidth}p{0.36\textwidth}p{0.32\textwidth}}
\toprule
攻击点 & 审稿人可能怎么说 & 论文里如何预防 \\
\midrule
\endhead
Construct = old idea renamed & TODO & 见 §[neighbor differentiation] \\
Levels not really 1D & TODO & 见 §[archetype framing] \\
Formative 倒推 & TODO & 见 §[alternative outcomes] \\
Ground-truth instability & TODO & 见 §[pairwise preference + disagreement] \\
Deployment cost hidden & TODO & 见 §[cost as finding] \\
Demand characteristic & TODO & 见 §[naming audit] \\
Generalization beyond testbed & TODO & 见 §[scope statement] \\
TODO ($\geq$ 5 rows for deep) & TODO & TODO \\
\bottomrule
\end{longtable}
\end{center}

% =====================================================================
% Block 12 — Demand Characteristic Audit
% =====================================================================

\subsection{命名与提示的预埋审计}

\begin{center}
\small
\begin{tabularx}{\textwidth}{>{\bfseries}p{0.16\textwidth}p{0.20\textwidth}p{0.20\textwidth}p{0.10\textwidth}X}
\toprule
Surface & Current wording & Embedded answer & Risk & Fix \\
\midrule
System name & TODO & TODO & TODO & TODO \\
Condition labels & TODO & TODO & TODO & TODO \\
Main RQ wording & TODO & TODO & TODO & TODO \\
Interview prompts & TODO & TODO & TODO & TODO \\
\bottomrule
\end{tabularx}
\end{center}

% =====================================================================
% Block 13 — Elevator Pitch Box
% =====================================================================

\section{一段话电梯演讲}

\begin{thesisbox}
本项目不再把重点放在 TODO（obvious framing），而是研究一个更基础的 TODO（field）问题：
TODO（sharpened question）。

在 TODO（real setting）, TODO（phenomenon explanation）。
因此 TODO（system / method）不能 TODO（naive goal）, 而要 TODO（precise goal）。
我把这个问题命名为 \textit{[construct]}。

通过 TODO（testbed）, 我会先 TODO（Study 1）, 再 TODO（Study 2/3）, 最后 TODO（Study 4）。
本论文的贡献不是 TODO（naive thing）, 而是 TODO（precise contribution）。
\end{thesisbox}
